\documentclass[12pt,letterpaper]{article}

\usepackage[margin=0.75in]{geometry}
\usepackage{array}
\usepackage{tabularx}
\usepackage[hidelinks]{hyperref}

\pagestyle{empty}
\setlength{\parindent}{0pt}
\setlength{\parskip}{0.45em}
\setlength{\emergencystretch}{2em}

\newcommand{\field}[1]{\textbf{#1}}
\newcolumntype{Y}{>{\centering\arraybackslash}X}

\begin{document}

\begin{center}
  \textbf{\large MATH 15200: Calculus II}\\
  \textbf{Autumn 2026}
\end{center}

\vspace{0.4em}
\hrule
\vspace{1.2em}

\field{Instructor:} Edward Varvak.\\
\field{Office:} Stevanovich 113.\\
\field{Email:} \href{mailto:varvak@uchicago.edu}{varvak@uchicago.edu}.\\
\field{Class schedule:} Tuesdays and Thursdays, 8:00--9:20 AM (Section 14) in Eckhart 308.\\
\field{Office hours:} Tuesdays and Thursdays, 9:30--10:30 AM in Stevanovich 113, and by appointment.

\vspace{0.5em}
\field{VCA/Grader:} Aryan Kaushal.\\
\field{Email:} \href{mailto:aryankaushal@uchicago.edu}{aryankaushal@uchicago.edu}.\\

\field{Textbook:} \textit{Calculus: One and Several Variables}, 10th edition, by Salas, Hille, and Etgen.\\

\field{Course content:} Integration of functions of one real variable, its applications and techniques, transcendental functions, and elementary differential equations. The course will begin with a review of key material from MATH 15100.\\

\field{Course description:} This course provides a detailed introduction to integral calculus, blending theoretical and computational work. Students should be able to state and understand the principal definitions and theorems, carry out computations, and solve applied problems. Short proofs may appear in lecture or on homework, but proofs are not the primary focus.\\

\field{Prerequisites:} MATH 15100 or equivalent. Students should be comfortable with limits, continuity, differentiation, the Mean Value Theorem, algebra, and trigonometry. We will review the formal definition of a limit at the start of the quarter.\\

\field{Exams:} Two in-class midterms are tentatively scheduled for Thursday, October 22, and Thursday, November 19. There will also be a cumulative final examination; its date and time will be announced when available. All examination dates are tentative and subject to change.\\

\field{Homework:} One assignment will be due each Monday at 11:59 PM. It will include routine online problems completed through Edfinity, accessed via Canvas, and written problems uploaded to Canvas. Collaboration is permitted and encouraged, but each student must write and submit their own solutions independently. Written solutions should be clear and complete.\\

\field{Grading:} The final course grade will be calculated as follows: 10\% homework, 25\% Midterm I, 25\% Midterm II, and 40\% final examination.\\

\field{Policies:} Late homework is not accepted, but the lowest weekly homework score will be dropped. Absence from an examination will result in a zero except in truly extenuating circumstances. A student facing such circumstances should contact the instructor as promptly as possible.\\

\field{Accommodations:} Students approved for academic accommodations by Student Disability Services (SDS) should follow SDS procedures and contact the instructor promptly so that approved accommodations can be implemented. SDS can be reached at \href{tel:+17737026000}{(773) 702-6000}, \href{mailto:disabilities@uchicago.edu}{disabilities@uchicago.edu}, or \href{https://disabilities.uchicago.edu/}{disabilities.uchicago.edu}.\\

\field{Academic integrity:} Students must follow the University of Chicago's standards for academic integrity. Submitted work must be the student's own, subject to the homework collaboration policy above, and outside assistance must be acknowledged when appropriate. Suspected violations may receive a zero and be referred to the appropriate University office. Consult the College's \href{https://college.uchicago.edu/student-services/academic-integrity-student-conduct}{Academic Integrity and Student Conduct policy} and ask the instructor about any uncertain form of assistance.\\


\field{Tentative schedule:} The following schedule is tentative and may be adjusted as the quarter progresses. Section numbers refer to \textit{Calculus: One and Several Variables}, 10th edition, by Salas, Hille, and Etgen.

\begin{center}
\footnotesize
\renewcommand{\arraystretch}{1.5}
\setlength{\tabcolsep}{5pt}
\begin{tabularx}{\textwidth}{|>{\centering\arraybackslash}p{0.16\textwidth}|Y|Y|}
  \hline
  & \textbf{Tuesday} & \textbf{Thursday} \\
  \hline 
  \textbf{Week 1} &
  MATH 15100 review and the formal &
  Definition of the integral; Riemann sums\\
  Sep. 28--Oct. 2 & definition of a limit & \S\S5.1--5.2 \\
  \hline
  \textbf{Week 2} &
  Fundamental Theorem of Calculus &
  Properties of integrals; indefinite integrals  \\
  Oct. 5--9 & \S\S5.3--5.4 & \S\S5.6, 5.8 \\
  \hline
  \textbf{Week 3} &
  Substitution &
  Area; Mean Value Theorem for integrals \\
  Oct. 12--16 & \S5.7 & \S\S5.5, 5.9, 6.1 \\
  \hline
  \textbf{Week 4} &
  Volumes; Midterm I review &
  \textbf{MIDTERM I} \\
  Oct. 19--23 & \S\S6.2--6.3 & \textbf{October 22} \\
  \hline
  \textbf{Week 5} &
  Inverse functions; logarithms &
  Exponential functions; arbitrary powers \\
  Oct. 26--30 & \S\S7.1--7.3 & \S\S7.4--7.5 \\
  \hline
  \textbf{Week 6} &
  Exponential growth and decay &
  Integration by parts \\
  Nov. 2--6 & \S\S7.6--7.7 & \S8.2 \\
  \hline
  \textbf{Week 7} &
  Powers of trigonometric functions &
  Trigonometric substitution \\
  Nov. 9--13 & \S8.3 & \S8.4 \\
  \hline
  \textbf{Week 8} &
  Partial fractions; Midterm II review &
  \textbf{MIDTERM II} \\
  Nov. 16--20 & \S8.5 & \textbf{November 19} \\
  \hline
  \textbf{Week 9} &
  \multicolumn{2}{c|}{\textbf{THANKSGIVING BREAK --- NO CLASSES}} \\
  Nov. 23--27 & \multicolumn{2}{c|}{} \\
  \hline
  \textbf{Week 10} &
  Elementary differential equations &
  Final-exam review \\
  Nov. 30--Dec. 4 & \S\S9.1--9.2 & \\
  \hline
\end{tabularx}
\end{center}

College final examinations are scheduled for December 8--11. The date and time of the MATH 15200 final examination will be announced when they become available.

\end{document}
